\section{ТЕОРЕТИЧЕСКИЕ АСПЕКТЫ ЭКОЛОГИИ ЧЕЛОВЕКА}

\subsection{Влияние микропластика на физиологию человека}
Одной из самых острых проблем современной экологии человека является повсеместное распространение микропластика. Микропластик представляет собой мелкие частицы синтетических полимеров размером менее 5 миллиметров. Они образуются как в результате целенаправленного промышленного производства (первичный микропластик, используемый в косметике и бытовой химии), так и вследствие механического, ультрафиолетового и химического разрушения крупного пластикового мусора (вторичный микропластик).

Исследования последних лет показывают, что микропластик проник практически во все экосистемы планеты. Частицы полимеров обнаруживаются в водопроводной и бутилированной воде, в соли, меде, пиве и морепродуктах. Для жителя современного мегаполиса основным путем проникновения микропластика в организм является алиментарный (через пищу и воду) и ингаляционный (вдыхание пластиковой пыли, образующейся при трении автомобильных шин и износе синтетических тканей). Особую опасность представляет разогрев пищи в пластиковых контейнерах в микроволновой печи. Экспериментально доказано, что при нагревании даже пищевой пластик (маркировки PP и PET) может выделять миллионы микро- и наночастиц в пищу.

Попадая в желудочно-кишечный тракт и дыхательную систему, наночастицы пластика способны преодолевать эпителиальные барьеры, проникать в кровоток и накапливаться в паренхиматозных органах (печень, почки), а также проникать через гематоэнцефалический и плацентарный барьеры. Полимерные частицы не обладают острой токсичностью, однако они вызывают хронический окислительный стресс, локальные воспалительные реакции и нарушают микробиом кишечника. Кроме того, поверхность микропластика обладает высокой адсорбционной способностью: она притягивает и концентрирует тяжелые металлы, стойкие органические загрязнители (например, полихлорбифенилы) и патогенные микроорганизмы, выступая своеобразным «троянским конем» для токсинов.

\subsection{Качество питьевой воды в условиях мегаполиса}
Вода является фундаментальным элементом, обеспечивающим нормальное протекание всех биохимических процессов в организме человека. В условиях урбанизации обеспечение населения чистой питьевой водой становится сложной технологической и экологической задачей. Водозабор для крупных городов осуществляется преимущественно из поверхностных водоисточников, которые подвержены интенсивному антропогенному загрязнению: сбросам промышленных стоков, поверхностному стоку с автомобильных дорог и смывам сельскохозяйственных удобрений.

Современные водоочистные станции используют многоступенчатые системы подготовки, включая коагуляцию, флокуляцию, фильтрование и обеззараживание. Однако основным методом дезинфекции до сих пор остается хлорирование. Несмотря на свою высокую бактерицидную эффективность, активный хлор вступает в реакцию с остаточными органическими веществами, образуя хлорорганические соединения (например, трихлорметан), обладающие доказанным канцерогенным и мутагенным действием. 

Второй, не менее значимой проблемой, является состояние распределительных водопроводных сетей. Высокий процент износа труб приводит к так называемому вторичному загрязнению. Вода, идеально чистая на выходе с водоочистной станции, проходя километры по старым ржавым трубам, насыщается солями железа, тяжелыми металлами и биологическими пленками. Регулярное употребление такой воды без дополнительной доочистки в домашних условиях (например, фильтрами кувшинного типа или системами обратного осмоса) приводит к постепенному накоплению токсикантов в организме, что негативно сказывается на работе выделительной и сердечно-сосудистой систем.

\subsection{Цифровая гигиена и физиология сна}
Развитие информационных технологий радикально изменило образ жизни человека. Сегодня экранное время среднестатистического горожанина составляет от 6 до 10 часов в сутки. Постоянное взаимодействие со смартфонами, планшетами и мониторами породило новое направление в экологии человека — цифровую гигиену. Это комплекс правил и норм, направленных на минимизацию негативного влияния цифровой среды на физическое и психоэмоциональное здоровье.

Ключевым аспектом цифровой гигиены является регуляция воздействия экранов на циркадные ритмы человека. Эволюционно биологические часы человека настроены на смену дня и ночи. Сетчатка глаза содержит специализированные светочувствительные ганглиозные клетки (ipRGC), которые наиболее восприимчивы к свету синего спектра (длина волны около 460-480 нм). В дневное время синий свет солнца блокирует выработку мелатонина (гормона сна) в шишковидной железе, обеспечивая состояние бодрствования и концентрации внимания.

Современные LED-экраны гаджетов излучают интенсивный свет именно в синем спектре. Использование смартфонов в вечернее и ночное время подает мозгу ложный сигнал о наступлении светового дня. В результате происходит резкое подавление секреции мелатонина. Это приводит к удлинению латентного периода засыпания, нарушению фаз глубокого сна и снижению общей продолжительности отдыха. Хронический дефицит качественного сна напрямую коррелирует со снижением когнитивных функций, ослаблением иммунитета, развитием депрессивных состояний и нарушением метаболизма.

Соблюдение правил цифровой гигиены, таких как активация «ночного режима» (фильтрации синего света) на устройствах после 19:00, отказ от использования гаджетов за 1-2 часа до сна и удаление электронных устройств из спальной зоны, является критически важным для поддержания здоровья в эпоху информационных технологий. Именно просвещение в этих вопросах стало ключевой задачей разрабатываемого проекта «ЭкоМедиа».
